\chapter{Overactive Bladder}

\section*{Definition}
Overactive bladder (OAB) is a syndrome characterized by urinary urgency, usually with increased daytime frequency and nocturia, with or without urgency urinary incontinence. It affects approximately 12-17 percent of the adult population.

\section*{What Happens in Your Body}
OAB results from involuntary contractions of the detrusor muscle during the bladder filling phase (detrusor overactivity). This can be neurogenic (arising from neurological disease) or idiopathic. The normal bladder stores urine at low pressure; in OAB, inappropriate detrusor contractions cause the sensation of urgency. Mechanisms include increased afferent signaling, decreased central inhibition, and changes in urothelial function.

\section*{Causes}
\begin{itemize}
\item Idiopathic (most common, no identifiable cause)
\item Neurological conditions: stroke, Parkinson disease, multiple sclerosis, spinal cord injury
\item Bladder outlet obstruction (BPH in men, pelvic organ prolapse in women)
\item Age-related changes in bladder function
\item Urinary tract infection (acute, reversible)
\item Bladder stones or tumors
\item Diuretic use or excessive caffeine/alcohol intake
\item Cognitive impairment (impaired signaling)
\end{itemize}

\section*{When to See a Doctor}
\begin{itemize}
\item Urinary urgency: sudden, compelling desire to urinate that is difficult to defer
\item Urinary frequency: urinating 8 or more times in 24 hours
\item Nocturia: waking up 2 or more times at night to urinate
\item Urgency urinary incontinence: leaking urine with urgency
\item Small-volume voids
\item Sensation of incomplete emptying
\end{itemize}

\section*{Related Symptoms}
\begin{itemize}
\item Urinary urgency and frequency
\item Nocturia (nighttime urination)
\item Incontinence (urge urinary incontinence)
\item Enlarged prostate (BPH)
\item Urinary tract infection
\item Interstitial cystitis (bladder pain syndrome)
\end{itemize}

\section*{Self-Care / Home Management}
\begin{itemize}
\item Bladder training: scheduled voiding every 2-3 hours
\item Pelvic floor exercises (Kegel exercises) to strengthen the pelvic floor
\item Avoid bladder irritants: caffeine, alcohol, spicy foods, acidic fruits
\item Maintain a healthy weight
\item Reduce fluid intake in the evening
\item Double voiding: empty the bladder, wait, and try again
\end{itemize}

\section*{How Your Doctor Will Evaluate You}
\begin{itemize}
\item Bladder diary (frequency, volume, urgency episodes)
\item Urinalysis and urine culture to rule out infection
\item Post-void residual measurement
\item Urodynamic studies for complex or refractory cases
\item Cystoscopy if hematuria or bladder pathology suspected
\end{itemize}

\section*{Treatment Options}
\begin{itemize}
\item Behavioral therapy: bladder training, pelvic floor exercises, fluid management
\item Anticholinergic medications: oxybutynin, solifenacin, tolterodine, darifenacin
\item Beta-3 adrenergic agonists: mirabegron, vibegron
\item Percutaneous tibial nerve stimulation (PTNS)
\item Sacral neuromodulation (InterStim therapy)
\item Botulinum toxin type A injection into the detrusor
\item Surgery: augmentation cystoplasty or urinary diversion (rare)
\end{itemize}

\section*{References}
\begin{thebibliography}{99}
\bibitem{overactive_bladder:oab1} Gormley EA, Lightner DJ, Faraday M, et al. "Diagnosis and treatment of overactive bladder (non-neurogenic) in adults: AUA/SUFU guideline." J Urol. 2012;188(6 Suppl):2455-2463.
\bibitem{overactive_bladder:oab2} Nambiar AK, Bosch R, Cruz F, et al. "EAU guidelines on assessment and nonsurgical management of urinary incontinence." Eur Urol. 2018;73(4):596-609.
\bibitem{overactive_bladder:oab3} Andersson KE. "Antimuscarinics for treatment of overactive bladder." Lancet Neurol. 2004;3(1):46-53.
\bibitem{overactive_bladder:oab4} Chapple CR, Khullar V, Gabriel Z, et al. "The effects of antimuscarinic treatments in overactive bladder: a systematic review and meta-analysis." Eur Urol. 2005;48(1):5-26.
\bibitem{overactive_bladder:oab5} Abrams P, Cardozo L, Fall M, et al. "The standardisation of terminology of lower urinary tract function." Neurourol Urodyn. 2002;21(2):167-178.
\bibitem{overactive_bladder:oab6} Corcos J, Przydacz M, Campeau L, et al. "CUA guideline on adult overactive bladder." Can Urol Assoc J. 2017;11(5):E142-E173.
\end{thebibliography}
